\subsection{IMP Language}
\subsubsection{The syntax}
The allowed characters:
\begin{code}{haskell}
    Letter = 'A' | . . . | 'Z' | 'a' | . . . | 'z'
    Digit = '0' | '1' | '2' | '3' | '4' | '5' | '6' | '7' | '8' | '9'
\end{code}

And the allowed tokens:
\begin{code}{haskell}
    Ident = Letter { Letter | Digit }*
    Numeral = Digit | Numeral Digit
    Var = Ident
\end{code}

Arithmetic and Boolean statements could be defined in Haskell as follows
\begin{code}{haskell}
    data Aexp = Bin Op Aexp Aexp | Var String | Num Integer
    data Op = Add | Sub | Mul
    data Bexp = Or Bexp Bexp | And Bexp Bexp | Not Bexp | Rel Rop Aexp Aexp
    data Rop = Eq | Neq | Le | Leq | Ge | Geq
\end{code}
with the abbreviations for \texttt{Eq} (=), \texttt{Neq} (\#), \texttt{Le} ($<$), \texttt{Leq} ($\leq$), \texttt{Ge} ($>$) and \texttt{Geq} ($\geq$)

Statements could be defined in Haskell as follows
\mint{haskell}+data Stm = Skip | Assign String Aexp | Seq Stm Stm | If Bexp Stm Stm | While Bexp Stm+

We use the following naming conventions for meta-variables:
\begin{tables}{ll}{Variable     & Type}
              $n$               & for numerals (\texttt{Numeral})            \\
              $x, y, z$         & for variables (\texttt{Var})               \\
              $e, e', e_1, e_2$ & for arithmetic expressions (\texttt{Aexp}) \\
              $b, b_1, b_2$     & for boolean expressions (\texttt{Bexp})    \\
              $s, s', s_1, s_2$ & for Statements (\texttt{Stm}) \\
\end{tables}
Meta-variables stand for arbitrary program variables, whereas program variables are concrete variables in a program.

To denote \bi{syntactic equality} of two variables or statements, we use $\equiv$
